\documentclass[10pt,a4paper]{amsart}
\usepackage[margin=19mm]{geometry}
\usepackage{amsmath,amssymb,mathtools}
\usepackage[scr=rsfs,cal=euler]{mathalfa}
\usepackage{xfrac}
\usepackage[unicode,hidelinks,bookmarks=false]{hyperref}
\newcommand{\ie}{i.e.\ }
\newcommand{\cf}{cf.\ }
\newcommand{\resp}{respectively\ }
\newcommand{\Subsection}{\subsection}
\providecommand{\nobreakdash}{\nobreak\hbox{-}}
\setlength{\emergencystretch}{2em}
\multlinegap0pt
\usepackage{bm}
\usepackage{bbm}
\usepackage{dsfont}
\usepackage{xspace}
\usepackage{cancel}
\usepackage{mathtools}
\usepackage{amsthm}
\makeatletter
\@ifundefined{theorem}{\newtheorem{theorem}{Theorem}}{}
\@ifundefined{corollary}{\newtheorem{corollary}[theorem]{Corollary}}{}
\@ifundefined{definition}{\newtheorem{definition}[theorem]{Definition}}{}
\@ifundefined{lemma}{\newtheorem{lemma}[theorem]{Lemma}}{}
\makeatother
\newcommand{\acl}{\operatorname{acl}}
\newcommand{\tp}{\operatorname{tp}}
\title[Topologically 1-based T-minimal Structures]{Topologically 1-based T-minimal Structures}
\author{Benjamin Castle, Assaf Hasson, Will Johnson}
\date{Source: arXiv:2508.18558v1 (CC BY 4.0)}
\begin{document}
\maketitle

\noindent\emph{Source and licence.} Topologically 1-based T-minimal Structures. Version arXiv:2508.18558v1 (CC BY 4.0), distributed under the Creative Commons Attribution 4.0 International licence, as recorded in the arXiv licence metadata for this exact version and shown on the abstract page https://arxiv.org/abs/2508.18558v1. Full rights notice: licenses/arxiv-2508.18558-CC-BY-4.0.txt.\par\smallskip

\noindent\textit{Editorial scope.} Counted unit: the Theorem whose source label is \texttt{T: gconf} in the article (source0.tex lines 1457--1459, source label \texttt{T: gconf}), delivered together with the article's own complete original proof of it (source0.tex lines 1461--1500, one \texttt{proof} environment of 40 lines). No mathematics is written, repaired or re-derived: statement and proof are the article's own text line for line. Required context retained uncounted: L: mu generic cont (lines 722--729); D: weak type def (lines 1291--1298); T: combinatorial group conf (lines 1377--1379); L: family of bijections (lines 1396--1409); C: gconf is additive (lines 1441--1443). Every definition, display and label of those lines is retained unchanged. Omitted from this package and not redistributed: 1--721, 730--1290, 1299--1376, 1380--1395, 1410--1440, 1444--1456, 1460--1460, 1501--2358. Nothing inside a retained excerpt is omitted or altered: every retained body is delivered exactly as the article's lines are written, so no omission receipt of any kind accompanies this package. Citations: the retained excerpts carry no citation command at all, so no bibliography entry of the article is transcribed and no reference is redistributed. Reference adaptation: none was needed and none was made; every reference inside the delivered unit resolves inside it, so no unresolved reference or citation is produced. Printed slips kept literal: no printed word, symbol, hypothesis or number of the counted statement or of its proof was altered. Layout and class: the article's own class is \texttt{amsart} with packages including amssymb, hyperref, graphicx, float, xy; the wrapper is the lane's pinned \texttt{amsart} editorial wrapper at 10pt on A4, which restates the theorem-like environments the retained text needs and adds the article's own macro definitions that the retained text needs (\texttt{\textbackslash acl}, \texttt{\textbackslash tp}), with extra wrapper packages (bm, bbm, dsfont, xspace, cancel, mathtools). The delivered numbering is the unit's own. Assets: no retained span contains a figure environment, a graphics-inclusion command, a PDF-page inclusion, a table or a TikZ picture; no image file is copied into this package, the package declares no asset dependency and no image file is redistributed with it. Licence: version arXiv:2508.18558v1 of this article is distributed under the Creative Commons Attribution 4.0 International licence, as recorded in the arXiv licence metadata for this exact version and shown on the abstract page https://arxiv.org/abs/2508.18558v1. A full rights notice is delivered at licenses/arxiv-2508.18558-CC-BY-4.0.txt. Characters: every retained excerpt is pure ASCII, so the delivered engine, XeLaTeX with its default Latin Modern text font, reports no missing character.
\begin{lemma}\label{L: mu generic cont}
    Let $a\in M^m$, $b\in M^n$, and $A$ a parameter set so that $(a,b)$ is additive over $A$. Then: 
    \begin{enumerate}
        \item The projection $\mu(ab/A)\rightarrow\mu(b/A)$ is open and surjective.
        \item If $a\in\acl(Ab)$, then the projection $\mu(ab/A)\rightarrow\mu(b/A)$ is a homeomorphism.
        \item In particular, if $a$ and $b$ are interalgebraic over $A$, then $\mu(ab/A)$ is the graph of a homeomorphism $\mu(a/A)\rightarrow\mu(b/A)$.
    \end{enumerate}
\end{lemma}

\begin{definition}\label{D: weak type def}
    Let $A$ be a parameter set, and let $(I,<,\mathcal X,R,Mor)$ be a groupoid spine. We say that $(I,<,\mathcal X,R,Mor)$ is \textit{$A$-type-definable} if $I$ is small, each $X_i$ is $A$-type-definable, and the morphisms can be represented by $A$-type-definable families. More precisely, this means that for each $(i,j)\in R$, one can choose $A$-type-definable sets $T_{ij}$ and $F_{ij}\subset X_i\times X_j\times T_{ij}$, so that:
    \begin{enumerate}
        \item Each fiber $(F_{ij})_t\subset X_i\times X_j$ for $t\in T_{ij}$ is the graph of a (necessarily unique) morphism $i\rightarrow j$, which we denote $f^{ij}_t:X_i\rightarrow X_j$. 
        \item For each morphism $f:i\rightarrow j$, there is $t\in T_{ij}$ with $f^{ij}_t=f$.
    \end{enumerate}
    In this case, we call $\{T_{ij}\}$ and $\{F_{ij}\}$ a \textit{type-definition of $(I,<,\mathcal X,R,Mor)$ over $A$}.
    \end{definition}

\begin{theorem}\label{T: combinatorial group conf}
    Let $A$ be a parameter set, and let $(I,<,\mathcal X,R,Mor)$ be an $A$-type-definable regular groupoid spine with $|I|\geq 3$. Then for any $i$, there is an $A$-type-definable group in $\mathcal M^{eq}$ acting regularly and type-definably on $X_i$. In particular, for any $e\in X_i$, there is an $Ae$-type-definable group structure on the set $X_i$, with $e$ as the identity element.
\end{theorem}

\begin{lemma}\label{L: family of bijections} Let $(a,b,t)$ be real tuples. Assume:
\begin{enumerate}
\item Any two of $a,b,t$ are symmetrically independent.
\item $a$ and $b$ are interalgebraic over $t$.
\item $(t,ab)$ is additive. 
\item $\tp(ab/t)$ is topologially 1-based.
\end{enumerate}
Then $\mu(abt/\emptyset)$ is the graph of a regular family of homeomorphisms $\mu(a/\emptyset)\rightarrow\mu(b/\emptyset)$. More precisely, there is a family $F$ of homeomorphisms $f:\mu(a/\emptyset)\rightarrow\mu(b/\emptyset)$ such that:
\begin{enumerate}
    \item[(i)] For each $t'\in\mu(t)$, the fiber $\mu(abt/\emptyset)_{t'}\subset\mu(ab/\emptyset)$ is the graph of a homeomorphism $f_{t'}\in F$.
    \item[(ii)] For each $f\in F$, there is $t'\in\mu(t/\emptyset)$ with $f_{t'}=f$.
    \item[(iii)] For each $a'\in\mu(a/\emptyset)$ and $b'\in\mu(b/\emptyset)$, there is exactly one $f\in F$ with $f(a')=b'$.
\end{enumerate}
\end{lemma}

    \begin{corollary}\label{C: gconf is additive}
        Let $(a,b,c,t,u)$ be a one-dimensional pre-group configuration. If $A,B\subset\{a,b,c,t,u\}$, then $(A,B)$ is additive.
    \end{corollary}

    \begin{theorem}\label{T: gconf}
        Let $(a,b,c,t,u)$ be a one-dimensional topological group configuration in $\mathcal M$. Then there is an $\mathcal M$-type-definable group structure on the set $\mu(a/\emptyset)$ with identity element $a$ (and similarly for $b$ and $c$).
    \end{theorem}

    \begin{proof}
        We will build an $\mathcal M$-type-definable regular groupoid spine $(I,<,\mathcal X,R,\operatorname{Mor})$ in the structure $\mathcal N$, where:
        \begin{itemize}
            \item $I=\{1,2,3\}$ with the usual order.
            \item $X_1=\mu(a/\emptyset)$, $X_2=\mu(b/\emptyset)$, and $X_3=\mu(c/\emptyset)$.
            \item $R=\{(i,j):i<j\}$.
        \end{itemize}
        The Theorem will then follow by Theorem \ref{T: combinatorial group conf}. Precisely, Theorem \ref{T: combinatorial group conf} will provide an $\mathcal Ma$-type-definable group structure on the set $X_1=\mu(a/\emptyset)$, with identity element $a$ (and similarly for $b$ and $c$). But since $a$, $b$, and $c$ are tuples from $\mathcal M$, these groups are actually $\mathcal M$-type-definable. 

        Our task, then, is to define the sets $\operatorname{Mor}(i,j)$ for $i<j$. We will do this by applying Lemma \ref{L: family of bijections} to each of $(a,b,t)$, $(b,c,u)$, and $(a,c,tu)$. We leave it to the reader to check the various independence and algebraicity requirements needed to apply Lemma \ref{L: family of bijections} in each of these cases; note that there is no need to check the additivity requirements, as Corollary \ref{C: gconf is additive}  gives them for free.

        So, in order to apply Lemma \ref{L: family of bijections}, the main thing to check is that each of $\tp(ab/t)$, $\tp(bc/u)$, and $\tp(ac/tu)$ is topologically 1-based. But indeed:
        
        \begin{itemize}
            \item By assumption $t\in\acl(ab)$, so $\dim(t/ab)=0$, and thus $\dim(t/\operatorname{germ}(ab/t))$ must also be 0. Thus $\tp(ab/t)$ is topologically 1-based.
            \item By the same reason as above, $\tp(bc/u)$ is also topologically 1-based.
            \item Finally, $\tp(ac/tu)$ is topologically 1-based because of the assumption that $(a,b,c,t,u)$ is a one-dimensional topological group configuration.
        \end{itemize}
        
        So, we may legally apply Lemma \ref{L: family of bijections} in each of the above instances. Denote the three resulting families of bijections by $\operatorname{Mor}(i,j)$ (in the obvious way, so that $\operatorname{Mor}(i,j)$ is the family of maps from $X_i$ to $X_j$).
        
        If we can show that the data defined so far gives a groupoid spine, then it will automatically be $\mathcal M$-type-definable and regular. That is:
        \begin{itemize}
            \item Each $X_i$ is an infinitesimal neighborhood, so is $\mathcal M$-type-definable.
            \item Similarly, each $\operatorname{Mor}(i,j)$ is represented as in Definition \ref{D: weak type def} by an $\mathcal M$-type-definable family -- precisely  by $\mu(abt/\emptyset)$, $\mu(bcu/\emptyset)$, and $\mu(actu/\emptyset)$, respectively.
            \item By Lemma \ref{L: family of bijections}, each $\operatorname{Mor}(i,j)$ has the required property for the definition of regularity.
        \end{itemize}
        Thus, our only remaining task is showing that $(I,<,\mathcal X,R,\operatorname{Mor})$ indeed gives a groupoid spine. And for this, the only thing to check is closure under composition. So let $f\in\operatorname{Mor}(1,2)$, and let $g\in\operatorname{Mor}(2,3)$. We want to show that $g\circ f\in\operatorname{Mor}(1,3)$.

        Since $f\in\operatorname{Mor}(1,2)$, there is $t'\in\mu(t/\emptyset)$ so that the fiber $\mu(abt/\emptyset)_{t'}$ is the graph of $f$. Similarly, since $g\in\operatorname{Mor}(2,3)$, there is $u'\in\mu(u/\emptyset)$ so that the fiber $\mu(bcu)_{u'}$ is the graph of $g$. We will show that the fiber $\mu(actu/\emptyset)_{t'u'}$ is the graph of $g\circ f$. Let $h:\mu(a/\emptyset)\rightarrow\mu(c/\emptyset)$ be the bijection whose graph is the fiber above $(t',u')$. So we want to show that $g\circ f=h$.

        To show this, let $a'\in\mu(a/\emptyset)$. Note that $abctu\in\acl(atu)$ -- thus $\dim(atu)=3$, and thus $\mu(atu/\emptyset)=\mu(a/\emptyset)\times\mu(t/\emptyset)\times\mu(u/\emptyset)$. Moreover, by Lemma \ref{L: mu generic cont}, the projection gives a homeomorphism $$\mu(abctu/\emptyset)\rightarrow\mu(atu/\emptyset)=\mu(a/\emptyset)\times\mu(t/\emptyset)\times\mu(u/\emptyset).$$ Since $a'\in\mu(a/\emptyset)$, $t'\in\mu(t/\emptyset)$, and $u'\in\mu(u/\emptyset)$, we conclude that $(a',t',u')\in\mu(atu/\emptyset)$, and thus there are $b',c'$ so that $(a',b',c',t',u')\in\mu(abctu/\emptyset)$. Thus:
        \begin{itemize}
            \item $(a',b',t')\in\mu(abt/\emptyset)$, and so $b'=f(a')$;
            \item Similarly, $(b',c',u')\in\mu(bcu/\emptyset)$, and so $c'=g(b')$;
            \item And finally, $(a',c',t',u')\in\mu(actu/\emptyset)$, and so $c'=h(a')$.
        \end{itemize}
        
        It now follows that $g\circ f(a')=h(a')$. Since $a'\in\mu(a/\emptyset)$ was arbitrary, it follows that $g\circ f=h$, and this proves the theorem.
        \end{proof}

\end{document}
